\section{Kamera és videókezelés}

Az alkalmazás szenzoradatok mellett videós tartalom is megjelenik. A videós funkciók a végrehajtási
képernyőhöz kapcsolódnak. A munkamenet elején a felhasználó kiválaszthatja,
hogy milyen nézetet szeretne használni: csak a bemutatóvideót, csak az előlapi kamera képét,
vagy a kettőt egyszerre, egymás mellett. Ez rugalmasabb használatot tesz lehetővé, mert a felhasználó
eldöntheti, hogy inkább a minta alapján szeretné követni a gyakorlatot,
vagy a saját mozgását szeretné látni végrehajtás közben.

A bemutatóvideó minden feladat előtt rövid előkészítő szakaszban jelenik meg.
Ez alatt a felhasználó még nem végez mérést, hanem megfigyelheti, hogyan kell a következő
gyakorlatot végrehajtani. A tényleges adatgyűjtés csak az előkészítő szakasz után indul el,
így a rögzített szenzoradatok már a gyakorlat végrehajtásához kapcsolódnak. A kamera-előnézet
célja nem önálló videóelemzés, hanem a mozgás felhasználó általi vizuális követésének támogatása.


A videók kezelése részben platformfüggő feladat, ezért az alkalmazásban különválik a közös felület
és az Android-specifikus megvalósítás. A közös kódrészben a videólejátszáshoz tartozó absztrakció
szerepel, míg az Androidhoz kötődő működés külön implementációban kap helyet.




